\section{Conclusion and outlook}
\label{sec:conclusion}

We organized Holographic Polar Arithmetic and Omega Theory into a single operational chain connecting computation, geometry, and observation.
The central move is definitional: time and space are internal resource costs of a scan--readout protocol.

\paragraph{Time.}
Time complexity becomes physical time in two equivalent guises: scan depth (iteration count of $\Theta$) and 1D compilation depth (nearest-neighbor circuit depth implementing a local PQCA step). The Omega lapse $\lapse(x)=\kappa_0/\kappa(x)$ makes time dilation an explicit computational slowdown induced by routing overhead.

\paragraph{Space.}
Space complexity becomes readout resolution: the prefix length and canonical digit depth required to stabilize a coordinate under projection readout, together with the active-region workspace capacity $|R|\log d$.

\paragraph{Tradeoffs.}
A Weyl-pair uncertainty relation provides a hard complementarity between scan-localization and readout-localization, and reversible-computation theory explains why projection readout forces time--space tradeoffs rather than optional engineering compromises.

\paragraph{Open-endedness.}
Universal QCAs support undecidable local reachability predicates, supplying a theorem-level boundary for ``open-ended history.'' Modeling observers as interactive machines with oracle-like resources is an interpretation-layer interface that keeps unitary microdynamics intact.

\paragraph{Testability.}
Discrete scan dynamics induces an exact lattice dispersion relation: it is approximately linear at low energy and deviates at high momentum, providing a falsifiable template for energy-dependent propagation.

\subsection*{Future directions}
Two technically sharp directions follow from this organization.

\begin{itemize}
    \item \textbf{Quantitative $(T(\epsilon),S(\epsilon))$ bounds.} Derive explicit upper/lower bounds for stabilization time and readout space at tolerance $\epsilon$, connecting discrepancy control (Denjoy--Koksma/Ostrowski bounds) to complexity lower-bound techniques.
    \item \textbf{Linking $\kappa$ to observables.} Establish a more direct quantitative bridge from routing overhead $\kappa(x)$ to effective gravitational observables (time-delay, redshift templates), with controlled approximations and explicit error budgets.
\end{itemize}
